% !TeX root = ../main.tex
\subsection{تمرین \secnum{۵.۷}، استخراج احتمال‌های گیت‌ها}
\label{sec:assignment-power-probability}

\textbf{خواسته.}
احتمال‌های جدول \secnum{۵.۱} کتاب را از تابع منطقی هر گیت به دست
آورید \cite[جدول \secnum{۵.۱}، صفحه ۱۸۷ و تمرین \secnum{۵.۷}،
صفحه ۲۰۹]{weste}.

\textbf{حل.}
برای سیگنال \(X\)، \(p_X=\Pr(X=1)\) و
\(\bar p_X=1-p_X=\Pr(X=0)\) را تعریف می‌کنیم. فرض می‌کنیم
ورودی‌ها در یک دوره مستقل‌اند، پس احتمال یک ترکیب ورودی از ضرب
احتمال‌های آن به دست می‌آید. \(Y\) خروجی گیت است.

\begin{itemize}
  \item \lr{AND2} فقط در حالت \((1,1)\) خروجی یک دارد، پس
    \(p_Y=p_Ap_B\). \lr{AND3} فقط در حالت \((1,1,1)\)
    خروجی یک دارد، پس \(p_Y=p_Ap_Bp_C\).
  \item \lr{NOR2} فقط با دو ورودی صفر خروجی یک دارد، بنابراین
    \(p_Y=\bar p_A\bar p_B\). \lr{OR2} خروجی آن را معکوس می‌کند،
    پس \(p_Y=1-\bar p_A\bar p_B\) است.
  \item \lr{NAND2} مکمل \lr{AND2} است و \(p_Y=1-p_Ap_B\).
  \item \lr{XOR2} وقتی ورودی‌ها متفاوت‌اند خروجی یک دارد. دو
    حالت \((1,0)\) و \((0,1)\) هم‌زمان رخ نمی‌دهند، پس
    احتمال‌هایشان با هم جمع می‌شوند و
    \(p_Y=p_A\bar p_B+\bar p_Ap_B\) به دست می‌آید.
\end{itemize}

ستون احتمال جدول کتاب با این روابط به دست می‌آید. اگر علاوه‌بر
استقلال ورودی‌ها در هر دوره، بردار ورودی در دوره‌های متوالی مستقل
و هم‌توزیع باشد و بین نمونه‌ها گذار اضافی رخ ندهد، احتمال گذار
صعودی خروجی از ضرب احتمال صفر در دوره فعلی و احتمال یک در دوره
بعد به دست می‌آید \cite[رابطه \secnum{۵.۱۴}، صفحه ۱۸۷]{weste}.
\[
  \alpha_Y=\Pr(Y_k=0,Y_{k+1}=1)
          =(1-p_Y)p_Y.
\]
\(Y_k\) مقدار خروجی در دوره \(k\) است. به‌عنوان کنترل، با فرض
احتمال \(1/2\) برای هر ورودی، نتیجه گیت‌ها به ترتیب جدول
\ref{tab:assignment-power-probabilities} است.

\begin{table}[htbp]
  \centering
  \caption{احتمال یک‌بودن و \lr{Activity Factor} خروجی در حالت
    ورودی‌های مستقل و هم‌احتمال، با استقلال بین دوره‌ها. ستون رابطه
    از جدول \secnum{۵.۱} کتاب استخراج شده است \cite{weste}.}
  \label{tab:assignment-power-probabilities}
  \begin{tabular}{rccc}
    \toprule
    گیت & رابطه \(p_Y\) & \(p_Y\) در مثال & \(\alpha_Y\) در مثال \\
    \midrule
    \lr{AND2}  & \(p_Ap_B\) & \(1/4\) & \(3/16\) \\
    \lr{AND3}  & \(p_Ap_Bp_C\) & \(1/8\) & \(7/64\) \\
    \lr{OR2}   & \(1-\bar p_A\bar p_B\) & \(3/4\) & \(3/16\) \\
    \lr{NAND2} & \(1-p_Ap_B\) & \(3/4\) & \(3/16\) \\
    \lr{NOR2}  & \(\bar p_A\bar p_B\) & \(1/4\) & \(3/16\) \\
    \lr{XOR2}  & \(p_A\bar p_B+\bar p_Ap_B\) & \(1/2\) & \(1/4\) \\
    \bottomrule
  \end{tabular}
\end{table}
